\begin{figure*}[t]
    \centering
    \includegraphics[width=1.0\linewidth]{graphs/overall.pdf}
    % \caption{Overview architecture of our proposed PoG.}
    \caption{Overview of the PoG architecture. Exploration: After initialization (detailed in Figure \ref{fig:initial}), the model retrieves entity paths from $\mathcal{G}_q$ through three exploration phases. Path Pruning: PoG applies a three-step beam search to prune paths after each exploration phase. Question Answering: The pruned paths are then evaluated for question answering. If these paths do not fully answer the question, the model explores deeper paths until $D_{max}$ is reached or moves on to the next exploration phase.}
    \label{fig:overall}
    \vspace{-1em}
\end{figure*}
% \vspace{-2mm}
\section{Preliminary}
\label{sec:2_preliminary}
Consider a Knowledge Graph (KG) $\mathcal{G(E,R,T)}$, where $\mathcal{E}$, $\mathcal{R}$ and  $\mathcal{T}$ represent the set of entities, relations, and knowledge triples, respectively. 
Each knowledge triple $T \in \mathcal{T}$ encapsulates the factual knowledge in $\mathcal{G}$, and is represented as $T = (e_{h}, r, e_{t})$, where $e_{h}, e_{t} \in \mathcal{E}$ and $r \in \mathcal{R}$.
% $\mathcal{T = \{((e_{h},r,e_{t}) | e_{h}, e_{t} \in \mathcal{E}, r\in \mathcal{R}\}$ is a set of knowledge triples contains the abundant factual knowledge in  $\mathcal{G}$.
% ,where $r$ is the relation connecting head and tail entities. 
Given an entity set $\mathcal{E_S\subseteq E}$, the induced subgraph of $\mathcal{E_S}$ is denoted as $\mathcal{S=(E_S,R_S,T_S)}$, where
$\mathcal{T}_S = \{ (e, r, e') \in \mathcal{T} \mid e, e' \in \mathcal{E}_S \}$, and 
$\mathcal{R}_S = \{ r \in \mathcal{R} \mid (e, r, e') \in \mathcal{T}_S \}.$
% $\mathcal{T}_S = \{ (e, r, e') \in \mathcal{T} \mid e, e' \in \mathcal{E}_S \}$, and 
% $\mathcal{R}_S = \{ r \in \mathcal{R} \mid \exists e, e' \in \mathcal{E}_S : (e, r, e') \in \mathcal{T} \}
% $
% Furthermore, we denote $\mathcal{D}(e), \mathcal{D}(r)$ as short textual desription sets of each entity $e\in \mathcal{E}$, and each relation $r\in \mathcal{R}$. 
Furthermore, we denote $\mathcal{D}(e)$ and $\mathcal{D}(r)$ as the sets of short textual descriptions for each entity $e \in \mathcal{E}$ and each relation $r \in \mathcal{R}$, respectively.
For example, the text description of
the entity ``m.0f8l9c'' is $\mathcal{D}$(``m.0f8l9c'')= ``France''.
For simplicity, in this paper, all entities and relations are referenced through their $\mathcal{D}$ representations and transformed into natural language.


% \label{sec:2_preliminary}






% \begin{definition}[Knowledge Triple] 
% Given a KG $\mathcal{G}$,
% the abundant factual knowledge within $\mathcal{G}$ is represented by a set of knowledge triples $\mathcal{T}$, denotes as 
% $\mathcal{G} = \mathcal{T} = \{(e_{head},r,e_{tail})| e_{head}, e_{tail} \in \mathcal{E}, r\in \mathcal{R}\}$, where $r$ is the relation connecting head and tail entities. 
% \end{definition}

\begin{definition}[Reasoning Path] 
Given a KG $\mathcal{G}$,
a reasoning path within $\mathcal{G}$ is defined as a connected sequence of knowledge triples, represented as: $path_{\mathcal{G}}(e_1, e_{l+1}) = \{ T_1, T_2, ..., T_l \} = \{(e_1,r_1,e_2),(e_2,r_2,e_3)$ $,...,(e_{l},r_l,e_{l+1})   \}$ , where $T_i \in \mathcal{T}$ denotes the $i$-th triple in the path and $l$ denotes the length of the path, i.e., $length(path_{\mathcal{G}}(e_1, e_{l+1})) = l$.

\end{definition}


\begin{example}

% Consider a reasoning path between the entities "University" and "Student" in KG. The reasoning path is given by:
Consider a reasoning path between "University" and "Student" in KG:
$
path_{\mathcal{G}}(\text{University}$, $ \text{Student}) $ $= \{(\text{University}$, $\text{employs}$, $\text{Professor})$, $ (\text{Professor}$, $\text{teaches}$, $\text{Course})$, $ (\text{Course}$, $\text{enrolled\_in}$, $\text{Student})\}
$, and can be visualized as: 
\[
\text{University} \xrightarrow{\text{employs}} \text{Professor} \xrightarrow{\text{teaches}} \text{Course} \xrightarrow{\text{enrolled\_in}} \text{Student}.
\]
\noindent It indicates that a ``University" employs a ``Professor," who teaches a ``Course," in which a "Student" is enrolled. The length of the path is 3.

\end{example}



For any entity $s$ and $t$ in $\mathcal{G}$, if there exists a reasoning path between $s$ and $t$, we say $s$ and $t$ can reach each other, denoted as $s \leftrightarrow t$. 
The distance between $s$ and $t$ in $\mathcal{G}$, denoted as $dist_\mathcal{G}(s,t)$, is the shortest reasoning path distance between $s$ and $t$. For the non-reachable vertices, their distance is infinite.
% , denoted by $dist_\mathcal{G}(s,t)$= $+\infty$.
Given a positive integer $h$, the $h$-hop neighbors of an entity $s$ in $\mathcal{G}$ is defined as $N_\mathcal{G}(s, h)=\{t\in \mathcal{E} |dist_\mathcal{G}(s,t)\leq h\}$. 


% \begin{definition}[Entity Path] 
% Given a KG $\mathcal{G}$, a list of entity $list_e = [e_1, e_2, e_3, ..., e_1| e_1, e_2, e_3,...,e_l \in \mathcal{E}]$,
% the entity path of $list_e$ within $\mathcal{G}$ is defined as a connected sequence of reasoning path, represented as: $path_{\mathcal{G}}(list_e) = \{ path_{\mathcal{G}}(e_1,e_2), path_{\mathcal{G}}(e2,e3), ..., path_{\mathcal{G}}(e_{l-1},e_{l}) \}=\{ (e_s,r,e_t)|(e_s,r,e_t)\in P\}$, where $p =\{ path_{\mathcal{G}}(e_1,e_2), path_{\mathcal{G}}(e2,e3), ..., path_{\mathcal{G}}(e_{l-1},e_{l}) \}$

% \end{definition}

\begin{definition}[Entity Path]
Given a KG $\mathcal{G}$ and a list of entities $list_e$ = [$e_1, e_2, e_3, \ldots, e_l$], the entity path of $list_e$ is defined as a connected sequence of reasoning paths, which is denoted as $path_{\mathcal{G}}(list_e)$
$= \{path_{\mathcal{G}}(e_1, e_2), $ $
path_{\mathcal{G}}(e_2, e_3), \ldots, path_{\mathcal{G}}(e_{l-1}, e_l) \}=\{(e_s, r, e_t)$ 
$| (e_s, r, e_t) \in path_{\mathcal{G}}(e_{i}, e_{i+1}) \land 1 \leq i < l\}$.
% = $\{(e_s, r, e_t) | (e_s, r, e_t) \in \mathcal{T}\}$.
% Each $path_{\mathcal{G}}(e_i, e_{i+1})$ denotes a specific reasoning path between consecutive entities $e_i$ and $e_{i+1}$ in the list.
\end{definition}


% \begin{problem}[Knowledge Graph Question Answering (KGQA)]
% KGQA is a typical reasoning task based on KGs. Given a natural language question $q$ and a KG $\mathcal{G}$, the task aims to design a function $f$ to predict answers $a \in Answer(q)$ based on knowledge from $\mathcal{G}$, i.e., $a = f (q, \mathcal{G})$.
Knowledge Graph Question Answering (KGQA)
is a fundamental reasoning task based on KGs. Given a natural language question $q$ and a KG $\mathcal{G}$, the objective is to devise a function $f$ that predicts answers $a \in Answer(q)$ utilizing knowledge encapsulated in $\mathcal{G}$, \ie $a = f(q, \mathcal{G})$.
% Following previous works \cite{sun2019pullnet, rogluo2023reasoning, tog1.0sun2023think, tog2.0ma2024think}, 
Consistent with previous research \cite{sun2019pullnet, rogluo2023reasoning, tog1.0sun2023think, tog2.0ma2024think},
we assume the topic entities $Topic(q)$ mentioned in $q$ and answer entities  $Answer(q)$ in ground truth are linked to the corresponding entities in $\mathcal{G}$, i.e., $Topic(q) \subseteq \mathcal{E} \text{ and }  Answer(q) \subseteq \mathcal{E}$.
% \end{problem}


% \myparagraph{Problem statement}
% Given a knowledge graph $\mathcal{G}$ and a natural language question $q$, \textbf{P}ath-\textbf{o}ver-\textbf{G}raph (\textbf{PoG}) aims to utilize LLMs to address the KGQA task by exploring $\mathcal{G}$.

% Knowledge Graph Question Answering task by exploring topic entity paths within $\mathcal{G}$ until the question is sufficiently answered or the exploration ends.


% \myparagraph{Problem statement}
% Given a knowledge graph $\mathcal{G}$, a natural language question $q$, 
% and two positive integers $D_{max}$ and $W_{max}$ as the maximum depth and width of exploration, 
% the \textbf{P}ath-\textbf{o}ver-\textbf{G}raph (\textbf{PoG}) aims to utilize LLMs to address the KGQA task by providing a set of final reasoning paths, denote as $Paths_f(q) = path_{\mathcal{G}_q}(Topic(q))$.
% % The set of final reasoning paths is the entity paths of the topic entity from $q$ within the question graph $\mathcal{G}_q$, denotes as $Final\_paths = path_{\mathcal{G}_q}(Topic(q))$. 
% Here, question graph $\mathcal{G}_q$ is defined as an induced subgraph of $\mathcal{G}$, constructed based on the entity set $\mathcal{E}_q \subseteq \mathcal{E}$, which is a set of $D_{max}$-hops neighbors of each topic entity from question, \ie $\mathcal{E}_q = \{N_{\mathcal{G}}(e_q, D_{max}) \mid e_q \in Topic(q)\}$.
% % $\mathcal{E}_q = \{e \mid \forall e \in \mathcal{E}, \exists e_q \in Topic(q) : \text{dist}_{\mathcal{G}}(e, e_q) \leq D_{max}\}$,
% Additionally, the number of paths in $Final\_paths$ is constrained by $W_{max}$, \ie $|Paths_f(q)| \leq W_{max}$, ensuring that the search and reasoning processes are both efficient and manageable.




